\documentclass[11pt]{article}
\usepackage[margin=1in]{geometry}
\usepackage{amsmath,amssymb,amsthm}

\newcommand{\binomV}[2]{\binom{#1}{#2}}

\title{A one-colour exact-cover test for the first remaining case of Problem 835}
\author{}
\date{}

\begin{document}
\maketitle

\section*{The first layer}

In the usual reduction of Problem 835, a colouring for the case \(k\) gives a
large set \(S(k-1,k,2k)\).  The known necessary conditions leave the first
serious case \(k=16\).  The first-layer derivation used here may be stated as
follows.

Let \(V\) be a set of 21 points.  Suppose that
\[
  D:\binomV{V}{5}\longrightarrow \mathbb F_{17}
\]
is an \(LS(4,5,21)\), and suppose that there are eleven functions
\[
  g_j:\binomV{V}{6}\longrightarrow \mathbb F_{17},
  \qquad j=1,\ldots,11,
\]
such that each \(g_j\) point-extends \(D\), and for every 6-set \(U\subset V\)
one has
\[
  \{g_1(U),\ldots,g_{11}(U)\}
  =
  \mathbb F_{17}\setminus
  \{D(U\setminus\{u\}):u\in U\}.
\]
An obstruction to this first layer would rule out the \(k=16\) case.

\section*{One fixed colour}

Fix \(a\in\mathbb F_{17}\).  Put
\[
  D_a=\{S\in\binomV{V}{5}:D(S)=a\},
  \qquad
  G_{j,a}=\{U\in\binomV{V}{6}:g_j(U)=a\}.
\]
Then \(D_a\) is an \(S(4,5,21)\), and the sets \(D_a,G_{1,a},\ldots,G_{11,a}\)
must satisfy the following exact-cover equations:
\[
  \sum_{\substack{S\supset T\\ |S|=5}} d_S=1
  \quad\text{for all }T\in\binomV{V}{4},
\]
\[
  d_S+\sum_{v\in V\setminus S}x_{j,S\cup\{v\}}=1
  \quad\text{for all }j=1,\ldots,11,\ S\in\binomV{V}{5},
\]
and
\[
  \sum_{j=1}^{11}x_{j,U}
  +\sum_{\substack{S\subset U\\ |S|=5}}d_S=1
  \quad\text{for all }U\in\binomV{V}{6}.
\]
Here \(d_S\) is the indicator of \(S\in D_a\), and \(x_{j,U}\) is the indicator
of \(U\in G_{j,a}\).

\section*{Size of the test}

There are
\[
  \binom{21}{5}=20349
\]
\(d\)-variables and
\[
  11\binom{21}{6}=596904
\]
\(x\)-variables.  Thus the exact-cover matrix has \(617253\) rows.  The
columns are the \(4\)-set equations, the \((j,S)\)-equations, and the
\(6\)-set equations, namely
\[
  \binom{21}{4}+11\binom{21}{5}+\binom{21}{6}
  =5985+223839+54264
  =284088.
\]
Each column has length \(17\).  Counting row incidences gives
\[
  \binom{21}{5}(5+11+16)+11\binom{21}{6}(6+1)
  =4829496.
\]

\section*{Computational target}

The useful target is not to enumerate all \(LS(4,5,21)\)'s.  It is first to
decide whether the one-colour exact-cover instance above is satisfiable.  An
UNSAT certificate for this instance would already contradict the assumed
first-layer object.  If it is satisfiable, the next useful target is to extract
structural information from its solutions, especially automorphisms, forced
block counts on small subsets, and any obstruction to assembling all seventeen
colours simultaneously.

\end{document}
